\item \textbf{Walleye Trouble.} Walleye is a fish that is native to the northern United States and Canada that is known for its excellent taste. It is especially popular in Minnesota where it is deemed their state fish. As a result it is often overfished in Minnesota lakes, especially Lake Mille Lacs. One way to measure the health of the population is through the average length of the fish  because fish grow their entire lives. The DNR considers the walleye population in a lake to be strong if the mean length of walleye in the lake is more than 17 inches. The DNR conducted a fish survey and netted 660 walleye and calculated an average length of 17.5 inches and a standard deviation of 4.68 inches. Is this sufficient evidence to say the walleye population in Lake Mille Lacs is strong?  

  \begin{enumerate}
  \item Conduct a hypothesis test of the above situation, using a significance level of $\alpha$=0.05.
    \begin{enumerate}
    \item Identify the hypotheses $H_0$ and $H_A$.
         \vspace{0.3cm}   
    \item Find the test statistic.


\vspace{0.3cm}
\item Using the table below, select the appropriate value for the p-value based on your alternative hypothesis.
\begin{table}[h]
\centering
\begin{tabular}{l|l}
	\hline
	P(T $\textless$ t) & 0.991 \\
	P(T $\textgreater$ t) & 0.009 \\
	P(T $\textgreater$ $|$t$|$) & 0.018 \\
	\hline
\end{tabular}
\end{table}
\\
    \item  Provide an interpretation of the p-value you chose from the above table within the context of the problem.
\item[]
    \item Make a decision about $H_{0}$. Circle the correct answer.
      \begin{itemize}
      \item Reject the null hypothesis. The p-value $p < \alpha = 0.05$.
      \item Reject the null hypothesis. The p-value  $p > \alpha = 0.05$.
      \item Fail to reject the null hypothesis. The p-value $p < \alpha = 0.05$.
      \item Fail to reject the null hypothesis. The p-value  $p > \alpha = 0.05$.
      \end{itemize}
\item[]
    \item   Based on your decision, write a conclusion in the context of the problem.
    \end{enumerate}
  \item  (\textbf{EXTRA CREDIT})  Using the same sample information that you had in part (a), compute a 95\% confidence interval for the population mean.
  \item (\textbf{EXTRA CREDIT}) Using your Confidence Interval from above, how could of you have reached the same conclusion as you did with the hypothesis test from part a?

  \end{enumerate}